% !TEX root = ../main.tex
\clearpage
\phantomsection
\section*{7 \textbf{Ethics Considerations}}
\label{sec:ethics-considerations}
\addcontentsline{toc}{section}{7 Ethics Considerations}

Although this research relies exclusively on publicly available blockchain data, we will observe the following ethical principles to ensure responsible conduct. Because the study uses non-personal, public ledger data and does not intervene with human participants, it will align with UNISA's research ethics principles on privacy, confidentiality, responsible data management, and prior ethics reviews where required.


Ethics clearance process: before starting the empirical data collection and analysis phase, the researcher will submit the proposal, data-management plan, and public-data justification to the relevant UNISA School ethics process. If the study is classified as exempt or low risk because it uses public pseudonymous blockchain data and no human-participant intervention, the exemption or clearance confirmation will be retained as part of the research record. No attempt will be made to deanonymize wallet owners, combine wallet addresses with off-chain personal data, or publish address-level findings in a way that could reasonably identify individuals.


\proposalSubheading{1. Privacy and Anonymity}

\begin{itemize}
\item[\textbullet] \textbf{Pseudonymity Respect}: Arbitrum One addresses are pseudonymous, not inherently linked to real-world identities. We will not attempt to deanonymize addresses or correlate them with personally identifying information.
\item[\textbullet] \textbf{No Sensitive Data Collection}: We will restrict data collection to on- chain events (approve, permit, transfer, GMX \textbf{PositionDecrease}) and avoid scraping or integrating off-chain data such as IP addresses or KYC records.
\end{itemize}
\proposalSubheading{2. Data Usage and Compliance}

\begin{itemize}
\item[\textbullet] \textbf{Open Data Sources}: All queries will target public datasets (e.g., Google BigQuery \texttt{goog\_blockchain\_arbitrum\_one\_us}) in accordance with their terms of service and rate limits.
\item[\textbullet] \textbf{No Malicious Interaction}: We will not broadcast transactions or otherwise interact with smart contracts in ways that could disrupt network operations or degrade performance.
\end{itemize}
\proposalSubheading{3. Fairness and Bias Mitigation}

\begin{itemize}
\item[\textbullet] \textbf{Algorithmic Fairness}: We will audit EndorseRank and AWP for un- intended biases---e.g., systematically undervaluing new entrants or low- volume participants---and discuss these limitations in our analysis.
\item[\textbullet] \textbf{Transparency}: All code, parameter choices, and evaluation metrics will be documented and published, enabling peer review and independent replication.
\end{itemize}
\proposalSubheading{4. Responsible Disclosure}

\begin{itemize}
\item[\textbullet] \textbf{Vulnerability Reporting}: If our analysis uncovers potential vulnerabilities in smart contracts (e.g., unexpected approval patterns that could be exploited), we will confidentially notify the responsible protocol teams before public disclosure.
\end{itemize}
\proposalSubheading{5. Social and Financial Impact}

\begin{itemize}
\item[\textbullet] \textbf{No Financial Advice}: Although EndorseRank may inform risk assessment, we will explicitly disclaim that it is not financial advice and highlight the need for human oversight in lending and trading decisions.
\item[\textbullet] \textbf{Avoiding Misuse}: We will discuss potential adversarial manipulation (e.g., Sybil attacks) and recommend safeguards so that our methodology is not misused to target or disadvantage specific users.
\end{itemize}
By adhering to these practices, we ensure that our research advances on-chain reputation modeling ethically, respects user privacy, and contributes responsibly to the DeFi ecosystem.
